\begin{definition}\label{BLR.basicVerifier}\lean{BLR.basicVerifier}\leanok
Let \(V_{\mathrm{BLR}(\F_{2}, n)}^{f}\) be the verifier with oracle access to \(f: \F_{2}^n \to \F_{2}\) defined as follows:
\begin{enumerate}
  \item Sample \(x, y \leftarrow \F_2^n\).
  \item Accept if and only if \(f(x) + f(y) = f(x+y)\).
\end{enumerate}
\end{definition}

\begin{theorem}\label{BLR_basic_completeness_ZMod2}\lean{BLR_basic_completeness_ZMod2}\uses{BLR.basicVerifier}\leanok
Let \(n > 0\). For every linear function \(f: \F_{2}^n \to \F_{2}\), we have
\[\Pr \left[V_{\mathrm{BLR}(\F_{2}, n)}^{f} = 1\right] = 1.\]
\end{theorem}
\begin{proof}\leanok
For every sampled pair \(x,y\), linearity gives
\(f(x+y)=f(x)+f(y)\). Hence the verifier accepts for every random choice.
\end{proof}

\begin{theorem}\label{BLR_basic_soundness_ZMod2}\lean{BLR_basic_soundness_ZMod2}\uses{BLR.basicVerifier}\leanok
Let \(n > 0\) and \(f: \F_{2}^n \to \F_{2}\).
\[\Pr \left[V_{\mathrm{BLR}(\F_{2}, n)}^{f} = 0\right] \geq \Delta(f, \mathcal{LIN}).\]
\end{theorem}
\begin{proof}\leanok
% Lecture 07 first proves a weaker BLR bound and then gives this sharper
% finite-field bound by Fourier analysis; the Lean statement uses the sharper
% version specialized to \(\F_2\).
This is the \(F=\F_2\) specialization of the BLR soundness theorem: a function
that is \(\Delta(f,\mathcal{LIN})\)-far from linear violates the additive
identity \(f(x)+f(y)=f(x+y)\) with probability at least that distance.
\end{proof}
